% HD-F1-T21. GPT-6 Astra Pro; /root; 22 November 2026 UTC.
\section[HD-F1-T21: spin alternating source exclusion]{HD-F1-T21: every spin alternating perfect layer is excluded at the original curve source}
\label{sec:f1-spin-alternating-source}

The remaining T2-OTHER coefficient source has no almost
simple projective image with socle $A_n$, $n\ge7$, and
actual perfect central layer isomorphic to the standard
double cover $2.A_n$. This holds for every original
coefficient field in the remaining gate, every scalar
lift and every multiplicity on restriction to the layer.
The proof uses the original matrices and an actual
$2.A_6$ subgroup; it imposes no spin-irreducibility
assumption on that restriction.

\paragraph{1. The actual layer and the small commutator subgroup.}
Let $G\subset GL_K(V)$ be the original faithful absolutely
irreducible coefficient image, $m=\dim_KV\ge27$, and
$Z=Z(G)$ its scalar centre. Put $\overline G=G/Z$.
In the stated case its socle is $S=A_n$ and its actual
perfect layer is
\[
 L=[P,P]\triangleleft G,\qquad
 P=\text{inverse image of }S,\qquad
 \iota:2.A_n\xrightarrow{\sim}L.
 \tag{F1T21.1}
\]
The map $L\to S$ has kernel $Z(L)=\{1,z\}$.
Since $z$ is a nonidentity scalar involution in the
faithful original representation, its actual matrix
is $-1_V$.

Use the standard identification $\operatorname{Aut}(A_n)=S_n$
for $n\ge7$. One precise published restatement is the
opening of the proof of Rizo--Schaeffer Fry--Vallejo,
\emph{Galois action on the principal block and cyclic
Sylow subgroups}, Lemma~3.3. Thus every nonscalar
$g\in G$ induces conjugation by a nonidentity
permutation $\sigma\in S_n$ on the actual quotient
$S$. Some three-cycle $\tau\in A_n$ satisfies
$\sigma\tau\sigma^{-1}\ne\tau$: the three-cycles
generate $A_n$, and an automorphism fixing all of
them is the identity. Choose either actual lift
$t\in L$ of $\tau$ and set
\[
                   h=[g,t]\in L.
 \tag{F1T21.2}
\]
Its image is the nonidentity permutation
$\sigma\tau\sigma^{-1}\tau^{-1}$, supported on at
most six letters. Hence $h$ is noncentral.

The standard double covers have the actual subgroup
maps $2.A_6\hookrightarrow2.A_n\hookrightarrow2.S_n$,
compatible with the inclusions on letters and with
their common central involution. The complete proof
of Guralnick--Tiep, \emph{A problem of Koll\'ar and
Larsen on finite linear groups and crepant resolutions},
J.~Eur.~Math.~Soc.~14 (2012), Lemma~2.23, p.~621,
uses precisely these maps and the same six-letter
commutator containment. Only that standard subgroup
construction is imported here; its character-ratio
bound is replaced by the exact eigenvalue computation
in paragraph~2.

Here is the compatibility with the given $G$, whose
central extension is retained. Lift $\sigma$ to the
standard $2.S_n$. Its conjugation on $2.A_n$, transported
by $\iota$, and conjugation by $g$ induce the same
automorphism of $S$. Their quotient on each element
is a central scalar in $\{1,z\}$. The resulting map
$L\to\{1,z\}$ is a homomorphism, and is trivial
because $L$ is perfect. Thus the conjugation actions
agree exactly. A permutation of the support letters,
lifted in $2.S_n$, consequently identifies an actual
subgroup $H\subset L$ containing $h$ with $2.A_6$.
This is an inclusion into the original $G$, with its
central involution mapped to $z$; it is not a replacement
of the original coefficient representation.

\paragraph{2. Exact eigenspaces of the six faithful $2.A_6$ characters.}
Use the ordinary character table of $2.A_6=SL_2(9)$
in the published \emph{ATLAS of Finite Groups}, table
$A_6$ and its double cover. Its exact machine-readable
copy is GAP Character Table Library~1.3.9,
\texttt{CharacterTable("2.A6")}, data file
\texttt{ctoalter.tbl}; its recorded origin is the
\emph{ATLAS}, with power maps for $2,3,5$.
The group has order $720$ and a unique central
involution. The six irreducible rows on which this
involution acts as minus one have degrees
\[
                         4,4,8,8,10,10.
 \tag{F1T21.3}
\]
For an element in class $c$ of order $o$, its
$\exp(2\pi i a/o)$-eigenspace multiplicity in a
character $\chi$ is exactly
\[
 \mu_a(\chi,c)=\frac1o\sum_{j=0}^{o-1}
           \chi(c^j)\exp(-2\pi i aj/o).
 \tag{F1T21.4}
\]
Indeed averaging the actual cyclic-group action with
these coefficients is its eigenprojection; its trace
is both sides of (F1T21.4). This supplies the map
from the published table values to the required
eigenspace dimensions.

The complete calculation, over all six rows and all
eleven noncentral classes, gives the following maxima.
Each entry is the largest multiplicity over both rows
of the indicated degree, all classes of the listed
orders, and every eigenvalue:
\[
\begin{array}{c|rrrr}
 \chi(1)&o=4&o\in\{3,6\}&o=8&o\in\{5,10\}\\\hline
 4&2&2&1&1\\
 8&4&3&2&2\\
 10&5&4&3&2
\end{array}
 \tag{F1T21.5}
\]
The saved exact GAP certificate computes all $66$
row/class pairs by (F1T21.4), checks that every
multiplicity is a nonnegative integer and that their
sum is the degree, and also checks all character-row
inner products, the degree-square sum and the
class-size sum. The full multiplicity lists and
dependency hashes are retained. This is an exact
calculation from published table data; it does not
claim a fresh construction of that character table.

Every entry in (F1T21.5) is at most half the degree.
The actual restriction $V_{\mathbf C}|_H$ is a sum
of these six representations with nonnegative integer
multiplicities: complete reducibility follows by
averaging, and its centre acts as $-1$ on every
constituent. Therefore, for the particular noncentral
$h$ in (F1T21.2),
\[
 \dim\ker(h-\lambda)\le\frac m2\quad(\lambda\in\mathbf C),
                 \qquad\operatorname{rank}(h-1)\ge\frac m2.
 \tag{F1T21.6}
\]
All multiplicities of the original restriction are
included in this sum.

\paragraph{3. The original inertia contradiction.}
For every scalar $\lambda$, the exact matrix identity
on $V_{\mathbf C}$ is
\[
 h-1=\bigl((g-\lambda)t-t(g-\lambda)\bigr)g^{-1}t^{-1}.
 \tag{F1T21.7}
\]
Its rank is at most twice that of $g-\lambda$.
Combining (F1T21.6)--(F1T21.7) gives
\[
            \operatorname{rank}(g-\lambda)\ge\frac m4
            \qquad(g\notin Z,\ \lambda\in\mathbf C).
 \tag{F1T21.8}
\]
A nonidentity scalar has moving rank $m$. Thus every
original nonidentity inertia matrix $g_j$ satisfies
the same moving-rank lower bound. The original
parabolic dimension and period-rank requirements are
\[
       \sum_{j=1}^b\operatorname{rank}(g_j-1)=2m+6,
                         \qquad b\ge12.
 \tag{F1T21.9}
\]
Equations (F1T21.8)--(F1T21.9) give
$2m+6\ge3m$, hence $m\le6$, contradicting $m\ge27$.
This proves the stated spin-layer exclusion.

\paragraph{Geometric maps and exact scope.}
The map used is the actual subgroup composite
$2.A_6\simeq H\hookrightarrow L\hookrightarrow G$
acting on the original $V_{\mathbf C}$, followed by
the explicit eigenprojections and commutator identity.
The connected cover, coefficient local system, rational
primitive Jacobian projector, six-dimensional corner
and restricted theta form are those of the original
source. They retain their identity maps, with
polarization multiplier one, on the Poincare
biextension, parabolic pairing, variational Hodge and
de Rham realizations, Grothendieck group and Adams
operations. No subgroup-cover detector is substituted.
The result concerns the specified actual double-cover
layer; the other central-layer cases are determined
only by their separately proved results.
